\clearpage\begingroup
\newcommand{\DioptrikCurrencyTwelve}{%
\begin{tikzpicture}[x=.125pt,y=-.125pt,baseline=-10pt]
\path[use as bounding box] (0,0) rectangle (90,97);
\path[fill=black,even odd rule] (61,6) -- (60,7) -- (58,7) -- (52,13) -- (52,15) -- (51,16) -- (51,19) -- (50,20) -- (49,20) -- (48,21) -- (47,21) -- (45,23) -- (42,23) -- (41,24) -- (40,24) -- (35,29) -- (35,39) -- (36,40) -- (36,42) -- (37,42) -- (39,44) -- (39,45) -- (40,45) -- (41,46) -- (45,46) -- (51,52) -- (51,57) -- (48,60) -- (45,60) -- (44,59) -- (41,59) -- (38,56) -- (36,56) -- (35,55) -- (35,54) -- (31,50) -- (31,49) -- (28,46) -- (26,46) -- (23,49) -- (21,49) -- (16,44) -- (14,44) -- (12,46) -- (10,46) -- (9,47) -- (9,48) -- (8,49) -- (7,49) -- (6,50) -- (6,55) -- (5,56) -- (5,58) -- (8,61) -- (9,61) -- (10,62) -- (11,62) -- (12,63) -- (12,64) -- (13,65) -- (13,66) -- (16,69) -- (18,69) -- (21,66) -- (22,66) -- (25,69) -- (30,69) -- (31,68) -- (31,65) -- (30,64) -- (30,62) -- (27,59) -- (28,58) -- (30,58) -- (32,60) -- (32,61) -- (37,61) -- (38,62) -- (39,62) -- (41,64) -- (44,64) -- (46,66) -- (49,66) -- (54,71) -- (54,75) -- (53,76) -- (52,76) -- (52,81) -- (48,85) -- (41,78) -- (38,78) -- (37,79) -- (36,79) -- (35,80) -- (35,81) -- (34,82) -- (33,82) -- (32,83) -- (32,84) -- (31,85) -- (31,87) -- (32,88) -- (41,88) -- (42,87) -- (44,87) -- (48,91) -- (51,91) -- (52,90) -- (53,90) -- (55,88) -- (57,88) -- (60,85) -- (60,79) -- (61,78) -- (61,75) -- (63,73) -- (63,66) -- (65,64) -- (72,64) -- (81,55) -- (81,52) -- (83,50) -- (83,49) -- (84,48) -- (84,43) -- (83,42) -- (83,39) -- (81,37) -- (78,37) -- (78,43) -- (79,44) -- (79,49) -- (78,50) -- (78,51) -- (77,52) -- (77,53) -- (74,56) -- (72,56) -- (71,57) -- (71,58) -- (70,59) -- (63,59) -- (60,56) -- (60,52) -- (61,51) -- (61,49) -- (67,43) -- (69,43) -- (73,39) -- (73,37) -- (76,34) -- (76,33) -- (77,32) -- (77,29) -- (76,28) -- (76,20) -- (75,20) -- (74,19) -- (74,17) -- (73,16) -- (73,13) -- (71,11) -- (71,10) -- (69,8) -- (68,8) -- (66,6) -- cycle
(61,26) -- (62,26) -- (64,28) -- (66,28) -- (69,31) -- (69,36) -- (68,37) -- (68,38) -- (67,39) -- (66,39) -- (63,42) -- (61,42) -- (60,41) -- (60,39) -- (58,37) -- (58,29) -- cycle
(47,26) -- (48,26) -- (51,29) -- (51,31) -- (49,33) -- (49,41) -- (47,43) -- (45,43) -- (40,38) -- (40,36) -- (38,34) -- (38,33) -- (43,28) -- (45,28) -- cycle
(62,12) -- (64,12) -- (69,17) -- (69,22) -- (67,24) -- (65,24) -- (64,23) -- (61,23) -- (57,19) -- (57,17) -- cycle;
\end{tikzpicture}}

\SourceLeaf{158}\HistoricalHead{153}{R}{GAUSS AN ENCKE.}
\begin{center}
[12.]\par\medskip
\textbf{\textsc{Gauss} an \textsc{Encke}.\quad Göttingen, 2. Januar 1840.}
\end{center}
\noindent . . . . . . . . . . . . . . . . . . . . .\par
\noindent\hspace*{\parindent}Ebenso würden Sie mich verpflichten, wenn Sie mir etwas ausführlichere Auskunft über \textsc{Plössls} Arbeiten, namentlich sogen[annte] dialytische F[ern-]Röhre geben könnten. Irre ich nicht, so habe ich vor einigen Jahren irgendwo bei einer Prüfung der letztern eine Bezugnahme auf Sie gelesen. Ich würde mich also über eine Belehrung freuen (besonders wenn bezüglich auf \textls[120]{grössere} Fernröhre, als in seinem Preiscatalog aufgeführt sind) über die wesentlichsten Umstände, die derjenige wissen müsste, der ähnliches zu acquiriren wünscht: 1) Leistung, etwa durch einige Specialia erläutert, 2) Dimensionen, 3) Zubehör (Stativ eventuell Mikrometer, eingetheilte Kreise), 4) Preis, 5) Zeit, wie lange warten? Können Sie mir derartige Aufschlüsse geben, entweder über \textsc{Plössl} oder auch über \textsc{März}, so werde ich Ihnen dankbar sein; im entgegengesetzten Fall sehen Sie die Frage als gar nicht geschehen an. Ich möchte gern ein anständiges Sehwerkzeug acquiriren; bei der Armseligkeit oder Nullität der für hiesige Sternwarte bereiten Mittel allenfalls aus meiner Tasche. Natürlich ist von einem solchen grossen Fernrohr wie das Ihrige gar keine Rede. Auch eine parallatische Aufstellung hat ohne Neuen Bau hier keinen Werth, aber Azimuthalkreis und Höhenkreis würde ich dabei wünschen (die übrigens nur die Minute sicher und bequem geben dürfen), da man dann (vermöge einer von mir berechneten Hülfstafel[*)])\begingroup\renewcommand{\thefootnote}{}\footnotetext{[*) Von dieser Hilfstafel ist auch in \textsc{Gauss}’ Briefen an \textsc{Schumacher} vom 7. und 12. September 1844 (\textit{Briefwechsel zwischen Gauss und Schumacher} IV, S. 286 und 292 ff.) die Rede. Vergl. die \textit{Tafel, um für eine bestimmte Polhöhe aus dem Stundenwinkel und der Declination, Azimuth, Höhe und parallactischen Winkel zu berechnen} in der \textit{Sammlung von Hülfstafeln}, herausgegeben 1822 von H. C. \textsc{Schumacher}, neu herausgegeben und vermehrt von G. H. L. \textsc{Warnstorff}, Altona 1845, S. 135; diese Tafel gilt für die Polhöhe von Altona und ist nach \textsc{Gauss}’ Angaben berechnet. Die in den Briefen an \textsc{Schumacher} beschriebene, von \textsc{Gauss} selbst schon im Jahre 1824 berechnete Tafel bezog sich auf die Polhöhe von Göttingen; siehe auch weiter unten unter \textit{Astronomie}.]}\endgroup\ \textls[120]{fast eben so schnell und bequem}, wie mit einer parallatischen Aufstellung, jeden Gegenstand finden kann. Können Sie mir also dabei rathen, indem der Maasstab etwa der sein müsste, dass die Kosten höchstens 1000\,\DioptrikCurrencyTwelve\ oder wenigstens nicht viel darüber betrügen? Ich bin etwas mistrauisch gegen öffentliche Anpreisungen geworden; \textsc{Sartorius} hat ein \textsc{Plössl}sches Fernrohr
\par\HistoricalFoot{XI1.}{20}
\clearpage\SourceLeaf{159}\HistoricalHead{154}{L}{PHYSIK.\quad BRIEFWECHSEL.}
\noindent erhalten, wovon er nicht viel rühmen will, ich weiss aber keine nähern Umstände. Ebenso hatte \textsc{Meyerstein} ein Objectiv von \textsc{März} erhalten, was ich gerne als Ablesefernrohr beim Magnetometer hätte gebrauchen mögen, da es stärkere Öffnungen und Vergrösserungen hat, als die jetzt von mir gebrauchten (die Vergrösserung war 60mahl, die sonstigen Dimensionen weiss ich nicht genau mehr). Allein mit der Wirkung dieses Objectives war ich durchaus nicht zufrieden. Dagegen habe ich einige (freilich nur kleinere terrestrische Zug-)Fernröhre von einem gewissen \textsc{Vogtländer} in Wien kürzlich in Händen gehabt, die ich in ihrer Art vortrefflich fand, und ganz entschieden besser, als \textsc{Frauenhofer}sche von derselben Grösse, die ich selbst besitze. . . . . . . .
\bigskip\begin{center}
[13.]\par\medskip
\textbf{\textsc{Encke} an \textsc{Gauss}.\quad Berlin, 12. Januar 1840.}
\end{center}
\noindent\{. . . . . . . . . . . . . . . . . . . . .\par
\noindent\hspace*{\parindent}Von \textsc{Plössls} Arbeiten kenne ich nur ein kleines zweifüssiges Fernrohr, was zu den ersten dialytischen gehört, die er gemacht hat, und was er als eines der besseren den Naturforschern in Wien präsentirte. Es ist nemlich eines der vier mit einem Sternchen bezeichneten. Herr v. \textsc{Buch} kaufte es damals für 140 fl. und hat es hier auf die Sternwarte gegeben. Dieses Fernrohr ist in der That vorzüglich. Seine Vergrösserung geht bis 80mal und entzückte auch \textsc{Hansen} als er hier war so sehr, dass er sogleich den Entschluss fasste, sich auch eines kommen zu lassen. \textsc{Frauenhofer}sche habe ich nicht damit vergleichen können, aber ich ziehe es den hiesigen guten englischen dreifüssigen bei weitem bei gleicher Vergrösserung vor. --- Grössere \textsc{Plössl} habe ich nicht gesehen. Es sind mehrere für die Navigationsschulen angeschafft worden, aber direct dahin abgegangen. . . . . . .\}
\clearpage\endgroup
